\documentclass[11pt,a4paper]{article}

\usepackage[T1]{fontenc}
\usepackage[utf8]{inputenc}
\usepackage[brazil]{babel}
\usepackage{geometry}
\usepackage{booktabs}
\usepackage{longtable}
\usepackage{tabularx}
\usepackage{hyperref}
\usepackage{enumitem}
\usepackage{amsmath}
\usepackage{microtype}
\usepackage{setspace}

\geometry{margin=2.2cm}
\setlist[itemize]{leftmargin=*}
\setlist[enumerate]{leftmargin=*}
\hypersetup{colorlinks=true,linkcolor=black,urlcolor=blue,citecolor=black}

\title{Pareamento probabilístico Censo\,$\times$\,CPF\\
\large Modelo com nome da mãe e endereço\\
\large Recorte UF\,21}
\author{Projeto probabilístico --- IBGE}
\date{Outubro de 2026}

\begin{document}
\maketitle
\tableofcontents
\newpage

%----------------------------------------------------------------------
\section{Introdução}
%----------------------------------------------------------------------

Este relatório descreve o pipeline de \emph{record linkage} entre o Censo
Demográfico~2022 e a base de CPF, no recorte da UF\,21 (Maranhão). O objetivo
operacional é atribuir, a cada pessoa do Censo elegível, no máximo um CPF, com
controle de falso positivo e cobertura aceitável.

A abordagem é probabilística (Splink, \texttt{link\_only}): o modelo estima a
probabilidade de match de cada par candidato e uma etapa posterior de
\emph{associação} transforma essas notas em atribuições com unicidade no
registro censitário e limite operacional de até três Censos por CPF.

Em relação à versão anterior do pipeline---em que a mãe entrava apenas como
prova na atribuição---esta versão inclui \textbf{nome da mãe} e
\textbf{endereço} (logradouro, CEP, município) diretamente no \emph{score} do
modelo, e remove a UF do score. Abaixo de $p = 0{,}90$, a atribuição exige
corroboradores explícitos (regras R1/R2/R3) para não repetir um piso solto de
nota.

A validação principal usa uma lista de ouro 1:1 (pares Censo--CPF conhecidos).
As métricas reportadas de erro entre associações e de não-associação são as da
\emph{associação} final, não as curvas pairwise do Splink.

%----------------------------------------------------------------------
\section{Visão geral das bases --- Brasil}
%----------------------------------------------------------------------

Antes do linkage, Censo e CPF são descritos de forma independente no bronze
nacional (\texttt{01c\_estatisticas\_descritivas\_bases}, sem filtro de UF).
O limpo tratado nacional ainda não foi materializado; a parametrização usou
a UF\,21 (seção seguinte).

\subsection{Volumes e cobertura geográfica}

\begin{center}
\begin{tabular}{@{}lrrp{5.8cm}@{}}
\toprule
Base & Registros & IDs distintos & Observação \\
\midrule
CPF bronze & $275{.}895{.}048$ & $275{.}895{.}048$ & 1 linha por CPF \\
CPF nomes fonéticos & $275{.}895{.}048$ & $275{.}895{.}048$ & coberta $100\%$ \\
Censo pessoas & $203{.}053{.}689$ & --- & questionário 2022 \\
Censo nomes & $195{.}101{.}203$ & $195{.}101{.}203$ & 1 por \texttt{ID\_MORADOR} \\
\bottomrule
\end{tabular}
\end{center}

As duas fontes cobrem as \textbf{27 UFs}. No CPF, $18{.}849{.}931$ registros
($6{,}8\%$) estão sem código de município (UF nula na agregação); o restante
distribui-se pelas UFs, com maiores volumes em SP ($52{,}4$\,M), MG
($26{,}5$\,M) e RJ ($22{,}8$\,M). O Maranhão (UF\,21) tem $8{.}447{.}170$ CPF
e $6{.}769{.}584$ pessoas no Censo bronze.

Gap Censo pessoas $\times$ nomes: $\sim$8\,M pessoas sem linha em
\texttt{censo\_pes\_nome} (não entram no limpo com nome).

\subsection{Completude no bronze}

\begin{center}
\small
\begin{tabular}{@{}lrr@{}}
\toprule
Campo & CPF bronze (\%) & Censo / nomes (\%) \\
\midrule
nome & $100{,}0$ & $100{,}0$ (tabela de nomes) \\
nome fonético & $100{,}0$ & $100{,}0$ \\
data de nascimento & $100{,}0$ & --- (no limpo UF\,21: $85{,}1\%$) \\
sexo & $100{,}0$ & --- \\
nome da mãe & $93{,}7$ & --- (no limpo UF\,21: $33{,}8\%$) \\
município & $93{,}2$ & $100{,}0$ (id da pessoa) \\
CEP & $100{,}0$ & --- \\
logradouro & $99{,}4$ & --- \\
idade calculada (Censo) & --- & $100{,}0$ \\
\bottomrule
\end{tabular}
\end{center}

Sexo no CPF bronze (código bruto): $51{,}18\%$ código~1, $48{,}80\%$ código~2,
$0{,}02\%$ código~9 ($58{.}383$ registros). Sem duplicidade de chave: $0$
linhas extras no CPF e em \texttt{censo\_pes\_nome}.

\subsection{Tamanho dos nomes (amostra aleatória de $\approx$2\,M por fonte)}

No bronze nacional, a distribuição de tokens do nome fonético foi estimada
com \textbf{amostra aleatória} de aproximadamente 2 milhões de registros de
cada fonte (\texttt{USING SAMPLE \ldots ROWS} no DuckDB, amostragem tipo
reservoir).

\begin{center}
\begin{tabular}{@{}lrr@{}}
\toprule
Tokens (fonético) & CPF (\%) & Censo nomes (\%) \\
\midrule
1 & $\sim 0$ & $0{,}5$ \\
2 & $7{,}4$ & $20{,}0$ \\
3 & $58{,}4$ & $52{,}2$ \\
4 & $31{,}4$ & $25{,}4$ \\
5+ & $2{,}8$ & $2{,}0$ \\
\bottomrule
\end{tabular}
\end{center}

No Censo nacional há bem mais nomes curtos (2 tokens: $20\%$ vs $7\%$ no CPF),
padrão que na UF\,21 também aparece entre quem fica sem par.

%----------------------------------------------------------------------
\section{Recorte experimental --- UF\,21 (Maranhão)}
%----------------------------------------------------------------------

A UF\,21 foi o ambiente para definir comparisons, blocking, faixas de
associação e validação na lista de ouro, antes de escalar ao Brasil.

\subsection{Universo após filtro e limpeza}

\begin{center}
\begin{tabular}{@{}lrrr@{}}
\toprule
Etapa & Censo & CPF & Nota \\
\midrule
Bronze filtrado UF\,21 & $6{.}775{.}805$ & $8{.}447{.}170$ &
  CPF de $275{.}895{.}048$ \\
Registros (staging) & $6{.}775{.}805$ & $8{.}447{.}170$ & \\
Limpo (\texttt{00b}) & $6{.}589{.}804$ & $8{.}447{.}170$ &
  $-186{.}001$ Censos sem nome \\
Aplicação (sem ouro exclusão) & $6{.}079{.}030$ & $7{.}931{.}465$ &
  tira lista determinística \\
\bottomrule
\end{tabular}
\end{center}

Óbito desligado (\texttt{ANO\_OBITO\_CORTE}$=0$). Chaves no limpo: $0$
duplicata de \texttt{unique\_id} / \texttt{cpf\_norm}.

\subsection{Completude no limpo (Censo $\times$ CPF)}

\begin{center}
\small
\begin{tabular}{@{}lrr@{}}
\toprule
Campo & Censo (\%) & CPF (\%) \\
\midrule
nome completo & $100{,}0$ & $100{,}0$ \\
data de nascimento & $85{,}1$ & $97{,}0$ \\
sexo & $100{,}0$ & $99{,}99$ \\
nome da mãe & $33{,}8$ & $96{,}9$ \\
município & $100{,}0$ & $100{,}0$ \\
CEP & $89{,}0$ & $100{,}0$ \\
logradouro & $100{,}0$ & $98{,}9$ \\
nome do meio & $88{,}4$ & $96{,}0$ \\
último nome & $99{,}6$ & $100{,}0$ \\
\bottomrule
\end{tabular}
\end{center}

No CPF limpo restam $1{.}156$ sexos nulos ($0{,}01\%$; código~$9$ na origem).
Diferença crítica para o pareamento: \textbf{mãe} ($34\%$ vs $97\%$) e data
($85\%$ vs $97\%$). Logradouro no Censo vem do CNEFE (quase sempre preenchido).

\subsection{Sexo, idade, tokens e municípios}

\begin{itemize}
  \item Sexo: Censo $50{,}85\%$ F / $49{,}15\%$ M; CPF $49{,}21\%$ F /
        $50{,}78\%$ M.
  \item Idade preenchida: CPF $97{,}0\%$ (média $37{,}2$, mediana $33$);
        Censo $14{,}9\%$ na coluna de idade (média $34{,}4$) --- o modelo
        usa sobretudo a data ($85\%$ com DOB).
  \item Primeiro nome mais frequente no limpo (\texttt{01}): \textbf{Maria}
        (Censo $579{.}880$; CPF $772{.}583$).
\end{itemize}

Tokens do nome fonético (limpo, universo completo):

\begin{center}
\begin{tabular}{@{}lrr@{}}
\toprule
Tokens & Censo (\%) & CPF (\%) \\
\midrule
1 & $0{,}4$ & $\sim 0$ \\
2 & $11{,}2$ & $4{,}0$ \\
3 & $51{,}9$ & $53{,}9$ \\
4 & $34{,}3$ & $39{,}4$ \\
5+ & $2{,}2$ & $2{,}7$ \\
\bottomrule
\end{tabular}
\end{center}

Maiores municípios (limpo): São Luís ($2111300$) --- $992{.}067$ Censo /
$1{.}493{.}884$ CPF; Imperatriz ($2105302$); São José de Ribamar ($2111201$).

\subsection{Conjunto de referência (lista de ouro)}

\begin{center}
\begin{tabular}{@{}lr@{}}
\toprule
Quantidade & $n$ \\
\midrule
Pares na lista de avaliação (nacional no arquivo) & $28{.}172{.}369$ \\
desses, 1:1 (Censo e CPF únicos) & $28{.}005{.}791$ \\
\textbf{G = 1:1 com ambos no limpo UF\,21} & $\mathbf{661{.}722}$ \\
Censos com CPF da lista de exclusão no limpo & $510{.}774$ \\
\bottomrule
\end{tabular}
\end{center}

Arquivos: exclusão da aplicação \texttt{regras1a2}; avaliação
\texttt{regras1a5} (V5). A ouro de avaliação é tipicamente mais fácil que o
Censo médio; recall/erro da validação valem em $G$, não se transportam cruos
aos $4{,}3$\,M atribuídos.

\subsection{Por que parametrizar na UF\,21}

\begin{enumerate}
  \item Volume compatível com treino EM, predict e validação interativa.
  \item Assimetría mãe/data já medida (tabela acima) --- stress-test das
        variáveis novas.
  \item Onomástica concentrada e mais nomes curtos no Censo, que pressionam
        blocking e desambiguação.
  \item Lista de referência no estado ($G=661{.}722$) para erro entre
        associações realizadas.
\end{enumerate}

A configuração (comparisons, 14 blockers, faixas $0{,}90$/$0{,}50$/$0{,}25$,
R1--R3) é candidata a reaplicação nacional --- ainda não medida no Brasil
inteiro.

%----------------------------------------------------------------------
\section{Pipeline e notebooks}
%----------------------------------------------------------------------

O fluxo é sequencial. Treino e aplicação são separados: o modelo é estimado no
limpo completo do recorte; a predição operacional roda no limpo \emph{sem} a
lista ouro determinística de exclusão.

\subsection{Preparação e diagnóstico}

\begin{description}[style=nextline]
  \item[\texttt{00\_preparar\_bases}]
    Importa bronze Censo + CPF, filtra UF/município, infere mãe, CEP e
    logradouro (CNEFE), carimba \texttt{cpf\_norm} da coorte no Censo e grava
    \texttt{censo\_registros} / \texttt{cpf\_registros} (sem empilhar as
    fontes).
  \item[\texttt{00b\_limpar\_dados}]
    Limpeza (óbito, nome obrigatório, sentinelas~$\to$~NULL) produz
    \texttt{censo\_limpo} / \texttt{cpf\_limpo}. Remove a lista ouro
    operacional da aplicação
    (\texttt{*\_limpo\_aplicacao}).
  \item[\texttt{01} / \texttt{01b}]
    Análise descritiva no limpo cheio (treino) e no limpo de aplicação.
\end{description}

\subsection{Modelo mãe + endereço}

\begin{description}[style=nextline]
  \item[\texttt{02c\_treinar\_splink\_mae\_endereco}]
    Define comparisons, estima prior, $u$, $m$ e EM; grava
    \texttt{splink\_model\_mae\_endereco.json}.
  \item[\texttt{02d\_aplicar\_splink\_mae\_endereco}]
    Carrega o JSON, aplica as 14 regras de blocking de predição e faz
    \texttt{predict} com $p \ge 0{,}25$ no limpo de aplicação; exporta
    \texttt{splink\_predictions\_mae\_endereco.parquet}.
\end{description}

\subsection{Atribuição e avaliação}

\begin{description}[style=nextline]
  \item[\texttt{04\_atribuir\_mae\_endereco}]
    Escada modelo / resgate / baixo (R1--R3) com teto cumulativo de 3 Censos
    por CPF; exporta \texttt{splink\_atribuicao\_mae\_endereco.parquet}.
  \item[\texttt{05\_avaliar\_mae\_endereco}]
    Cobertura, qualidade dos atribuídos, decomposição dos sem CPF e amostras
    por faixa/regra.
  \item[\texttt{06\_numeros\_finais}]
    Fatias no \texttt{censo\_limpo}: lista ouro + pareamento + sem CPF.
  \item[\texttt{03c\_avaliar\_lista\_ouro\_mae\_endereco}]
    Validação na lista de ouro de avaliação: mesmo modelo, blocking e
    atribuição do pipeline; pacote recall / precisão / FP / FN.
\end{description}

Notebooks da escada antiga (\texttt{02}, \texttt{02b}, \texttt{04},
\texttt{03b}, \texttt{04b}--\texttt{04e}) permanecem no repositório como
referência da versão anterior e não entram nos números deste relatório.

%----------------------------------------------------------------------
\section{Variáveis escolhidas}
%----------------------------------------------------------------------

O Splink compara quatro conceitos (comparisons). A UF \textbf{não} entra no
score nesta versão.

\subsection{Nome completo fonético (\texttt{nome\_completo\_phon})}

Níveis (do mais forte ao ELSE), com ajuste de frequência no match exato:

\begin{itemize}
  \item match exato;
  \item prefixo de tokens ($\ge 2$ tokens);
  \item Damerau--Levenshtein $\le 1$ no nome completo;
  \item um token diferente com DL $\le 2$ (último token igual);
  \item até dois tokens com DL $\le 2$;
  \item um token a mais/menos;
  \item mesmos tokens em ordem diferente;
  \item último token igual e Jaro--Winkler $\ge 0{,}95$ / $\ge 0{,}92$;
  \item ELSE.
\end{itemize}

\subsection{Data de nascimento (\texttt{data\_nascimento})}

Null conjunto (data e partes/idade); match exato com TF; mês/dia trocados;
dia diferente; mês diferente; ano $\pm 1$; DL $\le 1$ e $\le 2$ na data;
demais diferenças com datas presentes; fallback por idade exata / $|\Delta|=1$;
ELSE.

\subsection{Nome da mãe (\texttt{nome\_mae\_phon})}

Null; match exato com TF; nível \textbf{Contida} (tokens do nome mais curto
todos no mais longo, $\ge 2$ tokens); mesmos tokens em ordem diferente;
Jaro--Winkler $\ge 0{,}92$ e $\ge 0{,}85$; ELSE.

\subsection{Endereço composto (\texttt{endereco})}

Null / não comparável; logradouro+CEP exact; logradouro+município exact;
logradouro similar (JW~$\ge 0{,}90$) + CEP exact; CEP exact; ELSE.

Colunas auxiliares usadas em blocking e nas regras de associação (não são
comparisons próprias): \texttt{primeiro\_nome\_phon},
\texttt{ultimo\_nome\_phon}, partes da data, \texttt{sexo},
\texttt{cod\_municipio}, \texttt{cep}, \texttt{logradouro\_norm},
\texttt{idade}.

%----------------------------------------------------------------------
\section{Regras de treinamento}
%----------------------------------------------------------------------

Tipo de linkage: \texttt{link\_only} (Censo $\times$ CPF). Treino no limpo
completo.

\subsection{Prior e $u$}

\begin{itemize}
  \item Regras determinísticas do prior (distintas da predição):
        \texttt{cpf\_norm}; \texttt{nome\_completo\_phon} +
        \texttt{data\_nascimento}. Recall assumido $0{,}7$.
  \item $u$ por amostragem aleatória de pares
        (\texttt{U\_MAX\_PAIRS} $= 5{\cdot}10^8$).
\end{itemize}

\subsection{Estimativa de $m$ e EM}

\begin{enumerate}
  \item $m$ direto pela coluna \texttt{cpf\_norm}
        (\texttt{estimate\_m\_from\_label\_column}).
  \item EM em \texttt{data\_nascimento} + \texttt{sexo} (treina comparação de
        nome; $u$ fixo; até $10^7$ pares).
  \item EM em primeiro+último nome (treina comparação de data).
  \item EM em nome completo + data (treina comparação de mãe).
  \item EM em primeiro+último + data (treina comparação de endereço).
\end{enumerate}

Há ainda uma sessão diagnóstica com sexo no bloco de data, que \textbf{não}
fica no modelo salvo. O JSON final é
\texttt{splink\_model\_mae\_endereco.json}.

%----------------------------------------------------------------------
\section{Blocking de predição}
%----------------------------------------------------------------------

As 14 regras abaixo geram candidatos no \texttt{predict} (união OR). Não são
o blocking do prior nem do EM. Ficam de fora \texttt{nome\_meio} e
\texttt{cpf\_norm} na predição (quem não tem meio sumiria; CPF na predição
circularia a validação).

\begin{enumerate}[leftmargin=*,itemsep=0.15em]
  \item \texttt{nome\_completo\_phon}
  \item primeiro + último + ano
  \item primeiro + último + mês + dia
  \item primeiro + último + mês + ano
  \item primeiro + data completa
  \item último + data completa
  \item primeiro + mês + dia + município
  \item primeiro + mês + ano + município
  \item último + mês + dia + sexo + CEP
  \item último + mês + ano + sexo + CEP
  \item último + ano + mês + sexo + primeiro ``parecido''
        (Damerau--Levenshtein proporcional, comprimento $\ge 6$)
  \item primeiro + último + município + sexo + $|\Delta$idade$| \le 1$
  \item logradouro + CEP + município + sexo + $|\Delta$idade$| \le 1$
  \item mãe fonética + município + sexo + $|\Delta$idade$| \le 1$
\end{enumerate}

Na aplicação (UF\,21), o predict grava pares com $p \ge 0{,}25$.

%----------------------------------------------------------------------
\section{Regras de associação}
%----------------------------------------------------------------------

A associação transforma pares pontuados em atribuições. Um Censo recebe no
máximo um CPF; um CPF recebe no máximo três Censos (teto cumulativo entre
faixas). Em cada faixa, escolhe-se o melhor par elegível por $p$ e desempate
(mãe preenchida, logradouro, CEP).

\subsection{Faixas}

\begin{center}
\begin{tabular}{@{}llp{7.2cm}@{}}
\toprule
Faixa & Nota $p$ & Decisão \\
\midrule
modelo & $p \ge 0{,}90$ & aceita pelo score \\
resgate & $0{,}50 \le p < 0{,}90$ & R1 ou R2 ou R3 \\
baixo & $0{,}25 \le p < 0{,}50$ & R1 forte ou R2 (\textbf{sem R3}) \\
--- & $p < 0{,}25$ & não resgata \\
\bottomrule
\end{tabular}
\end{center}

\subsection{Corroboradores}

\begin{description}
  \item[R1]
    Data exata + primeiro nome compatível (igual ou Levenshtein $\le 1$) +
    (nome contido \textbf{ou} cobertura de tokens $\ge 0{,}75$; no baixo,
    cobertura $\ge 0{,}80$).
  \item[R2]
    Nome completo igual/contido + data plausível distinta da exata:
    $|\Delta$dias$| \le 1$, ou mesmo dia/mês com $|\Delta$ano$| \in \{1,10\}$.
  \item[R3]
    Primeiro compatível + cobertura nominal $\ge 0{,}55$ + mãe presente +
    Jaro--Winkler da mãe $\ge 0{,}75$. Somente no resgate
    ($0{,}50$--$0{,}90$).
\end{description}

A remoção de R3 do baixo veio de amostras com nome forte, mãe parecida e data
distante (risco clássico de falso positivo).

%----------------------------------------------------------------------
\section{Resultados}
%----------------------------------------------------------------------

\subsection{Cobertura na aplicação (UF\,21)}

Universo de aplicação: $6{.}079{.}030$ Censos (já sem a lista ouro
determinística de exclusão).

\begin{center}
\begin{tabular}{@{}lrr@{}}
\toprule
Versão & Atribuídos & Cobertura \\
\midrule
Escada antiga (mãe só na atribuição) & $4{.}092{.}388$ & $67{,}3\%$ \\
Piso simples $0{,}50$ (sem corroborador) & $4{.}352{.}929$ & $71{,}6\%$ \\
\textbf{Modelo + resgate + baixo (esta versão)} & $\mathbf{4{.}304{.}328}$ & $\mathbf{70{,}8\%}$ \\
\bottomrule
\end{tabular}
\end{center}

No \texttt{censo\_limpo}: lista ouro $7{,}8\%$, pareamento $65{,}3\%$, sem CPF
$26{,}9\%$ ($nos\_dois = 0$).

\subsubsection{Origem das atribuições}

\begin{center}
\begin{tabular}{@{}llrr@{}}
\toprule
Faixa & Regra & $n$ & \% dos atribuídos \\
\midrule
modelo & score & $4{.}130{.}020$ & $96{,}0\%$ \\
resgate & R1 & $62{.}354$ & $1{,}4\%$ \\
resgate & R2 & $20{.}897$ & $0{,}5\%$ \\
resgate & R3 & $20{.}040$ & $0{,}5\%$ \\
baixo & R1 & $62{.}882$ & $1{,}5\%$ \\
baixo & R2 & $8{.}135$ & $0{,}2\%$ \\
\bottomrule
\end{tabular}
\end{center}

Resgate+baixo = $174{.}308$ ($4{,}0\%$). Quase tudo é modelo alto.

\subsubsection{Sem CPF (melhor $p$)}

\begin{center}
\begin{tabular}{@{}lrr@{}}
\toprule
Melhor $p$ & $n$ & \% aplicação \\
\midrule
$\ge 0{,}50$ & $119{.}281$ & $2{,}0\%$ \\
$0{,}25$--$0{,}50$ & $132{.}252$ & $2{,}2\%$ \\
$0{,}05$--$0{,}25$ & $242{.}856$ & $4{,}0\%$ \\
sem par & $1{.}280{.}313$ & $21{,}1\%$ \\
\bottomrule
\end{tabular}
\end{center}

Dos $\sim$119\,k com $p \ge 0{,}50$ sem CPF, apenas 386 caem no teto; o resto
falhou R1/R2/R3. O buraco dominante é blocking ($1{,}28$\,M sem par).

\subsection{Falsos positivos e falsos negativos (lista de ouro)}

Validação no notebook \texttt{03c}: $G = 661{.}722$ pares 1:1 no limpo.
Métricas = associação final do pipeline completo (modelo + resgate + baixo).

\subsubsection{Definições}

\begin{itemize}
  \item \textbf{Acerto} (\texttt{associou\_cpf\_ouro}): associação com o CPF da
        verdade.
  \item \textbf{FP da associação}: associou outro CPF
        ($\mathrm{FP} = 1 - \mathrm{precisão}$ entre quem associou).
  \item \textbf{FN}: não associou o CPF ouro ($1 - \mathrm{recall}$); inclui
        quem associou errado, quem teve par abaixo da elegibilidade, e quem
        não entrou no predict.
  \item \textbf{Recall do par no predict}: o par ouro aparece no parquet com
        $p \ge 0{,}25$ (blocking + score, sem atribuição).
\end{itemize}

\subsubsection{Pacote principal}

\begin{center}
\begin{tabular}{@{}lrr@{}}
\toprule
Métrica & v1 (escada antiga) & v2 (mãe+endereço) \\
\midrule
Recall associação & $86{,}49\%$ & $\mathbf{91{,}67\%}$ \\
FN & $13{,}51\%$ & $\mathbf{8{,}33\%}$ \\
Precisão associação & $98{,}89\%$ & $\mathbf{99{,}37\%}$ \\
FP associação & $1{,}11\%$ & $\mathbf{0{,}63\%}$ \\
FP da escolha (antes do teto) & --- & $0{,}63\%$ \\
Recall do par ouro no predict & $92{,}03\%$ & $94{,}18\%$ \\
\bottomrule
\end{tabular}
\end{center}

Na ouro, o teto não retirou ninguém (FP da escolha = FP da associação).
Associações: $610{.}410$ (modelo $590{.}519$ + resgate $9{.}315$ + baixo
$10{.}576$).

\subsubsection{Decomposição do FN}

\begin{center}
\begin{tabular}{@{}lrr@{}}
\toprule
Corte & $n$ & \% da ouro \\
\midrule
Associou CPF ouro (acerto) & $606{.}585$ & $91{,}67\%$ \\
Associou CPF errado (FP e FN) & $3{.}825$ & $0{,}58\%$ \\
Teve par, não elegível / abaixo & $14{.}473$ & $2{,}19\%$ \\
Nenhum par no predict & $36{.}839$ & $5{,}57\%$ \\
\bottomrule
\end{tabular}
\end{center}

Leitura do FN ($55{.}137$ pessoas = $8{,}33\%$):
\begin{itemize}
  \item $67\%$ blocking vazio (\texttt{nenhum\_par});
  \item $26\%$ candidato existiu mas não passou na escada;
  \item $7\%$ associação ao CPF errado.
\end{itemize}

\subsubsection{Quadrantes nome $\times$ data (par ouro)}

\begin{center}
\begin{tabular}{@{}lrrr@{}}
\toprule
Recorte & $n$ & Recall & FP associação \\
\midrule
nome $\neq$, DOB $=$ & $152{.}038$ & $73{,}44\%$ & $1{,}12\%$ \\
nome $=$, DOB $\neq$ & $38{.}509$ & $67{,}03\%$ & $1{,}93\%$ \\
nome $=$, DOB $=$ & $471{.}175$ & $99{,}56\%$ & $0{,}44\%$ \\
\bottomrule
\end{tabular}
\end{center}

O ganho vs.\ v1 concentra-se em nome diferente com data igual (antes
$\sim 54\%$ de recall). O caso fácil permanece $\sim 99{,}5\%$.

\subsubsection{Nota sobre curvas Splink}

ROC / Precision--Recall com rótulos clericais medem ranking pairwise, não a
associação 1:1. Nesta rodada há poucos rivais no predict frente aos positivos
ouro; a curva PR pode parecer atípica. Os números operacionais são os do
pacote de associação acima.

%----------------------------------------------------------------------
\section{Limitações e o teto de cobertura ($\sim$71\%)}
%----------------------------------------------------------------------

Na aplicação, $70{,}8\%$ dos Censos receberam CPF. Os $29{,}2\%$ restantes
($1{.}774{.}702$ pessoas) \textbf{não} são um único fenômeno: a maior parte
nunca vira candidato no pareamento; uma fatia menor chega ao score e é
segurada de propósito pela associação.

\subsection{Decomposição dos sem CPF}

Entre quem ficou sem CPF, o melhor par no parquet de predição se parte assim:

\begin{center}
\begin{tabular}{@{}lrrp{6.2cm}@{}}
\toprule
Situação & $n$ & \% aplicação & Leitura \\
\midrule
Sem par no predict & $1{.}280{.}313$ & $21{,}1\%$ & blocking / piso do predict \\
Melhor $p$ em $[0{,}05;\,0{,}25)$ & $242{.}856$ & $4{,}0\%$ & candidato fraco \\
Melhor $p$ em $[0{,}25;\,0{,}50)$ & $132{.}252$ & $2{,}2\%$ & nota média, sem R1/R2 \\
Melhor $p \ge 0{,}50$ & $119{.}281$ & $2{,}0\%$ & nota boa, falhou R1/R2/R3
  (só 386 no teto de 3) \\
\bottomrule
\end{tabular}
\end{center}

Ou seja: \textbf{$\sim$3 em cada 4 sem CPF} estão no bloco
``sem par'' ($1{,}28$\,M / $1{,}77$\,M). Afrouxar o piso do modelo ou as
regras R1--R3 move no máximo os $\sim$8\% da aplicação que já têm candidato
com alguma nota; não mexe no buraco de $21$\,p.p.

\subsection{Perfil comparativo: atribuídos $\times$ sem par $\times$ sem CPF com par}

Diagnóstico no notebook \texttt{05}, sobre o Censo de aplicação.
Primeiro nome fonético mais frequente na UF: \textbf{MARIA} ($517{.}158$ pessoas).

\begin{center}
\small
\begin{tabular}{@{}lrrr@{}}
\toprule
Indicador & atribuído & sem par & sem CPF \emph{com} par \\
\midrule
$n$ & $4{.}304{.}328$ & $1{.}280{.}313$ & $494{.}389$ \\
sem data (\%) & $7{,}6$ & $\mathbf{36{,}2}$ & $38{,}5$ \\
sem mãe (\%) & $58{,}4$ & $72{,}3$ & $\mathbf{90{,}4}$ \\
sem CEP (\%) & $9{,}9$ & $14{,}4$ & $13{,}5$ \\
1º nome no top~10 (\%) & $24{,}9$ & $20{,}3$ & $\mathbf{32{,}2}$ \\
1º nome no top~50 (\%) & $38{,}0$ & $31{,}7$ & $44{,}0$ \\
nome com 2 tokens (\%) & $5{,}8$ & $\mathbf{30{,}0}$ & --- \\
nome com 3 tokens (\%) & $52{,}7$ & $44{,}7$ & --- \\
nome com 4 tokens (\%) & $38{,}9$ & $21{,}8$ & --- \\
\bottomrule
\end{tabular}
\end{center}

\begin{enumerate}
  \item \textbf{Sem data} é o maior contraste do sem par ($36\%$ vs $8\%$ nos
        atribuídos): várias das 14 regras de blocking usam ano/mês/dia/data.
  \item \textbf{Nome curto} ($\le 2$ tokens) concentra no sem par
        ($\sim 32\%$ vs $\sim 6\%$): menos âncora para regras
        primeiro+último+$\ldots$
  \item \textbf{Sem mãe} sobe no sem par ($72\%$) e é ainda pior em quem
        \emph{já tem} candidato ($90\%$): falta âncora para fechar a
        atribuição (e a 14ª regra de blocking).
  \item \textbf{Maria / top~10} não explicam o miss de blocking --- o sem par
        tem \emph{menos} top~10 que o atribuído. Nome frequentíssimo
        (Maria sozinha $\sim 517$\,k) torna a \textbf{desambiguação} difícil
        quando o candidato existe; por isso o pico de top~10 está no grupo
        ``sem CPF com par''.
  \item Cerca de $17\%$ do sem par tem data, nome, mãe, CEP e logradouro
        preenchidos: buraco fonético / sem correspondente no CPF, não missing.
\end{enumerate}

\subsection{Por que $1{,}28$\,milhões não entram no pareamento}

``Sem par'' significa: para aquele Censo, o \texttt{predict} não devolveu
nenhuma linha no parquet (nas 14 regras OR, contra o CPF limpo do estado).
Na prática isso junta dois efeitos, com o primeiro dominando:

\begin{enumerate}
  \item \textbf{Recall do blocking.}
    As regras exigem coincidência em combinações de nome fonético (completo,
    primeiro, último), partes da data, sexo, município, CEP, logradouro ou
    mãe. Quem tem nome muito divergente entre Censo e CPF, data ausente ou
    inconsistente, primeiro/último vazios, ou endereço/mãe sem preenchimento
    útil, simplesmente \emph{não gera candidato}---o modelo nunca vê o par.
  \item \textbf{Piso do predict.}
    Pares que só existiriam com nota muito baixa podem nem ser materializados,
    conforme o limiar exportado. Isso é secundário frente ao volume sem
    qualquer regra de bloco.
\end{enumerate}

Motivos estruturais por trás do miss de blocking (não são ``bug'' da escada):

\begin{itemize}
  \item qualidade e cobertura dos campos (mãe $\sim$34\% no Censo vs.\
        $\sim$97\% no CPF; logradouro/CEP com missings e normalização
        difícil);
  \item variação ortográfica/fonética além do que as 14 regras capturam;
  \item pessoas do Censo sem correspondente real ou encontrável no CPF do
        estado (óbito já filtrado à montante, migração, base incompleta,
        identificadores ruins);
  \item regras geográficas (município/CEP) que não disparam quando o endereço
        diverge ou está vazio---e aí a regra de mãe também não salva se a mãe
        estiver ausente no Censo.
\end{itemize}

Na lista de ouro o quadro é mais fácil: o recall do \emph{par} no predict já
é $94{,}2\%$ (só $5{,}6\%$ sem par). A ouro é uma população com par verdadeiro
1:1 conhecido; a aplicação inclui o caso médio e o caso difícil. Por isso a
ouro não ``promete'' 94\% de cobertura operacional.

\subsection{Por que a cobertura estabiliza perto de 71\%}

O teto observado não é arbitrário:

\begin{enumerate}
  \item \textbf{Teto do pareamento.}
    Mesmo um piso simples $p \ge 0{,}50$ \emph{sem} corroborador chegou a
    $71{,}6\%$---apenas $0{,}8$\,p.p.\ acima desta versão ($70{,}8\%$). Ou seja,
    com o blocking e o modelo atuais, há pouco a ganhar em cobertura sem
    mudar a geração de candidatos.
  \item \textbf{Escolha anti-FP.}
    Dos $\sim$119\,k com melhor $p \ge 0{,}50$ ainda sem CPF, quase todos
    falharam R1/R2/R3 (teto de 3 explica só 386). Essa fatia ($2{,}0\%$ da
    aplicação) é o preço explícito de não aceitar nota média/alta sem âncora.
    Na ouro, o mesmo desenho entrega FP de associação $0{,}63\%$ (vs.\
    $1{,}11\%$ na v1).
  \item \textbf{Faixa intermediária.}
    Outros $\sim$375\,k ($6{,}2\%$ da aplicação) têm algum par com
    $0{,}05 \le p < 0{,}50$: o modelo já duvida. Empurrá-los para atribuição
    sobe cobertura e muda o tipo de erro (homônimo / data frágil).
\end{enumerate}

Em uma conta fechada da aplicação:

\[
\begin{aligned}
70{,}8\% &\quad\text{atribuídos (modelo + resgate + baixo)}\\
+\ 2{,}0\% &\quad p\ge 0{,}50\text{ sem regra (reserva anti-FP)}\\
+\ 2{,}2\% &\quad 0{,}25\le p<0{,}50\\
+\ 4{,}0\% &\quad 0{,}05\le p<0{,}25\\
+\ 21{,}1\% &\quad\text{sem par (limite de blocking)}\\
&= 100\%.
\end{aligned}
\]

\subsection{O que esta versão não resolve (e o que resolveria)}

\begin{itemize}
  \item \textbf{Não resolve} o $1{,}28$\,M sem afrouxar ou redesenhar blocking
        (novas chaves, tolerância fonética maior, uso mais agressivo de mãe/
        endereço quando preenchidos, ou um segundo estágio de candidatos).
  \item \textbf{Não deve} ser lido como falha da atribuição: 96\% do que foi
        atribuído veio de $p \ge 0{,}90$, com FP baixo na ouro.
  \item \textbf{Caminhos para subir cobertura} com risco controlado: melhorar
        recall do blocking (prioridade); enriquecer mãe/endereço no Censo;
        só depois revisar R1--R3 ou o piso $0{,}90$ com validação ouro em
        mãos.
  \item \textbf{Limite externo:} o linkage não cria CPF inexistente nem
        corrige registro irrecuperável---parte do gap de $\sim$29\% é
        cobertura de bases e qualidade de dado, não só método.
\end{itemize}

%----------------------------------------------------------------------
\section{Conclusão}
%----------------------------------------------------------------------

O pipeline com mãe e endereço no score, piso de modelo $p \ge 0{,}90$ e
resgate/baixo com corroboradores R1/R2/(R3 só no meio) é a versão recomendada
para a UF\,21:

\begin{itemize}
  \item cobertura de aplicação $70{,}8\%$ ($+3{,}5$\,p.p.\ vs.\ escada antiga);
  \item na lista de ouro, recall sobe ($86{,}5\% \to 91{,}7\%$) e FP da
        associação cai ($1{,}11\% \to 0{,}63\%$);
  \item $96\%$ das atribuições vêm do modelo alto; o resgate é fino e
        ancorado;
  \item o teto perto de 71\% é sobretudo \textbf{blocking}
        ($1{,}28$\,M / $21\%$ sem par): a atribuição só explica $\sim$2\,p.p.\
        de reserva anti-FP; o piso $0{,}50$ solto já mostrava que não há
        muito mais cobertura barata com os candidatos atuais.
\end{itemize}

\end{document}
